Baseline fixed-effects estimates show that a 1 percent increase in AI adoption is associated with approximately 0.025 percent higher TFP (p-value below 0.001), controlling for human capital. The estimated AI-TFP relationship remains directionally stable under stronger identification checks, including two-way fixed effects (0.018, p-value 0.004) and trimmed-sample estimation (0.023, p-value below 0.001). For macro growth outcomes, the AI coefficient in the GDP growth specification is 0.004 (p-value 0.023), suggesting the productivity channel may emerge before fully translating into contemporaneous GDP growth. Sensitivity checks also support the main findings: the lagged AI specification yields an AI coefficient of 0.014 (p-value 0.047), and inference remains similar under time-clustered standard errors (0.018, p-value 0.009). The placebo outcome test using log human capital reports 0.011 (p-value 0.147), which should be interpreted as a falsification check rather than a causal estimate. A two-way fixed-effects specification with Driscoll-Kraay standard errors, which is robust to broad cross-sectional dependence, also preserves a positive AI coefficient (0.018, p-value 0.002).